% Ch.28 — QMTech XC7A100T FPGA (Hardware-Numerics Tier-1)
% φ² + φ⁻² = 3 · Trinity S³AI
% Author: Dmitrii Vasilev (ORCID 0009-0008-4294-6159)
% Issue: #422

\chapter{QMTech XC7A100T/200T FPGA: Hardware φ-Numerics Platform}
\label{ch:28}
\chaptermark{Ch.28 · QMTech FPGA}

% Header block (R7 compliance)
\begin{tcolorbox}[colback=gray!5,colframe=gray!40,title=Chapter Anchor]
\textbf{Anchor:} \(\varphi^2 + \varphi^{-2} = 3\) \\
\textbf{Theorems in chapter:} 4 (THM-28.1..28.4) \\
\textbf{Coq status:} \texttt{docs/phd/theorems/ch28\_fpga.v} \\
\textbf{FPGA artefact:} App.~\ref{app:fpga-bitstream}
\end{tcolorbox}

\section{Introduction}

This chapter documents the physical FPGA platform used to validate the
Trinity S\textsuperscript{3}AI φ-numeric coprocessor primitives. The platform
consists of a QMTech XC7A Starter Kit connected via a custom Xilinx Virtual
Cable (XVC) WiFi bridge to the development workstation.

The central result: the Period-Locked Monitor and φ-arithmetic cores
synthesise to 83~LUT / 27~FF — 0.06\% of the available resources — and
configure correctly on first attempt, with STAT register \texttt{0x401079FC}
confirming zero errors.

The anchor identity \(\varphi^2 + \varphi^{-2} = 3\) governs the period
ratios used in the monitor (Chapter~\ref{ch:24}) and the GoldenFloat encoding
scheme (Chapter~\ref{ch:6}).

\section{Device Identification}

\begin{theorem}[IDCODE Consistency — THM-28.1]
\label{thm:28:1}
The JTAG IDCODE \texttt{0x13631093} is a valid Xilinx Artix-7 IDCODE with
manufacturer \texttt{0x049}, part \texttt{0x3631} (XC7A200T), version~1.
All Trinity S\textsuperscript{3}AI design constraints are satisfiable on this
device.
\end{theorem}

\begin{proof}
IEEE~1149.1 IDCODE format: bit~0 = 1 (required), bits~1--11 = manufacturer,
bits~12--27 = part, bits~28--31 = version.
Parsing \texttt{0x13631093}:
\begin{align*}
\text{bit}~0 &= 1 \quad \checkmark \\
\text{manufacturer} &= (\texttt{0x13631093} \gg 1) \;\&\; \texttt{0x7FF} = \texttt{0x049} \quad (\text{Xilinx}) \\
\text{part} &= (\texttt{0x13631093} \gg 12) \;\&\; \texttt{0xFFFF} = \texttt{0x3631} \quad (\text{XC7A200T}) \\
\text{version} &= (\texttt{0x13631093} \gg 28) \;\&\; \texttt{0xF} = \texttt{0x1}
\end{align*}
XC7A200T contains a superset of XC7A100T resources.
The design requires 83~LUT $\ll$ 134,600~LUT (100T) $\ll$ 269,200~LUT (200T).
\end{proof}

\section{XVC WiFi Bridge}
\label{sec:xvc-bridge}

Standard JTAG cables (Digilent DLC10) are incompatible with macOS Sonoma
due to kernel-level USB driver conflicts (Appendix~J, BLK-001). The solution
is a custom XVC server running on an ESP32 microcontroller.

\subsection{Architecture}

\begin{verbatim}
macOS workstation
  └── openFPGALoader (XVC client)
       └── TCP 192.168.1.30:2542 (WiFi)
            └── ESP32 (XVC server)
                 └── GPIO: TMS=18, TCK=19, TDI=23, TDO=35
                      └── JTAG header P2 (QMTech board)
                           └── XC7A200T JTAG TAP
\end{verbatim}

\subsection{Key Implementation Detail}

The XVC \texttt{shift:} command handler was rewritten to process bits
individually rather than in 32-bit words, eliminating a systematic 4-bit
error at positions 14, 15, 16, 28 (BLK-003). The root cause was endianness
misalignment in the 32-bit packing path.

\begin{theorem}[XVC Correctness — THM-28.2]
\label{thm:28:2}
Bit-serial XVC processing produces IDCODE equal to the IEEE~1149.1 value
for all Artix-7 devices with IR length 6 and DR length 32.
\end{theorem}

\section{Resource Utilisation}

\begin{table}[h]
\centering
\begin{tabular}{lrrr}
\toprule
\textbf{Resource} & \textbf{Used} & \textbf{Available (100T)} & \textbf{Utilisation} \\
\midrule
LUT & 83 & 134,600 & 0.06\% \\
FF & 27 & 269,200 & 0.01\% \\
BRAM & 0 & 270 & 0\% \\
DSP & 0 & 240 & 0\% \\
IO & 6 & 285 & 2.1\% \\
\bottomrule
\end{tabular}
\caption{Post-route utilisation. Design is deliberately minimal — the FPGA
platform validates φ-arithmetic correctness, not throughput.}
\end{table}

\section{Verification Result}

\begin{theorem}[Configuration Success — THM-28.3]
\label{thm:28:3}
A bitstream synthesised targeting XC7A100T configures correctly on XC7A200T
if and only if (a)~no device-specific primitives are used, and (b)~all
constrained pins exist on the target package.
\end{theorem}

\begin{proof}
Condition (a): design uses only LUT6, FDRE, IBUF, OBUF — all portable across
the Artix-7 family. Condition (b): pins U18, D20, E19, R14, P14, N16, M16
are present in FGG484. Both conditions satisfied.
\end{proof}

The STAT register dump confirms all subsystems nominal:

\begin{verbatim}
STAT = 0x401079FC
DONE=1  EOS=1  GWE=1  MMCM_LOCK=1
DCI_MATCH=1  CRC_ERR=0  ID_ERR=0  DEC_ERR=0
\end{verbatim}

\begin{theorem}[Energy Baseline — THM-28.4]
\label{thm:28:4}
The FPGA idle power for 83~LUT / 27~FF at 50~MHz on XC7A200T is
$P_{\text{idle}} \leq 0.5$~W (static) + $P_{\text{dyn}} \leq 0.1$~mW
(dynamic at 0.06\% utilisation).
This establishes the hardware baseline for the energy comparison in
Chapter~\ref{ch:34}.
\end{theorem}

\section{Summary}

The QMTech XC7A200T platform (marketed as 100T) provides a production-grade
FPGA environment for Trinity S\textsuperscript{3}AI hardware validation. The
custom XVC WiFi bridge solves the macOS JTAG incompatibility (BLK-001), and
the bit-serial shift implementation resolves the sampling artefact (BLK-003).
All five blockers documented in Appendix~J are resolved. The platform is ready
for VSA hardware acceleration (Chapter~\ref{ch:31}) and energy benchmarking
(Chapter~\ref{ch:34}).

\(\varphi^2 + \varphi^{-2} = 3\)
